\chapter{Theory}\label{ch:theory}

\section{Finite-strain hyperelasticity}

A material point $\bm X$ of the reference body $\Omega_0$ moves to
$\bm x=\bm\varphi(\bm X)$. With the deformation gradient
$\F=\partial\bm\varphi/\partial\bm X$, $J=\det\F>0$ and $\Cten=\F^{\mathsf T}\F$,
a hyperelastic material is defined by a stored-energy density $\Psi(\F)$, from
which stress and tangent follow~\cite{Holzapfel2000}:
\begin{equation}
  P_{iJ}=\frac{\partial\Psi}{\partial F_{iJ}},\qquad
  \Aten_{iJkL}=\frac{\partial^2\Psi}{\partial F_{iJ}\,\partial F_{kL}}.
  \label{eq:PA}
\end{equation}
Because $\F$ and $\Pstress$ are not symmetric, the implementation keeps all nine
components and stores the tangent row-major as
$A^{9\times9}_{3i+J,\;3k+L}=\Aten_{iJkL}$. The proposal's $6\times6$ Voigt map
applies to symmetric strain--stress pairs and was therefore replaced. In the
total-Lagrangian setting, equilibrium is the stationarity of the potential,
\begin{equation}
  \int_{\Omega_0}\Pstress(\F):\nabla_{\!X}\delta\bm u\dd V=\delta W_{\mathrm{ext}}
  \quad\forall\,\delta\bm u,
  \label{eq:weak}
\end{equation}
and Newton's method needs $\Aten$ at exactly the same quadrature points as
$\Pstress$. The Phase~1 reference is the compressible neo-Hookean law
\begin{equation}
  \Psi_{\mathrm{NH}}=\tfrac{\mu}{2}(I_1-3)-\mu\ln J+\tfrac{\lambda}{2}(\ln J)^2,
  \qquad \mu=1,\ \lambda=10.
  \label{eq:nh}
\end{equation}

\section{Polyconvex neural energy}

A stored energy is \emph{polyconvex} if it is a convex function of
$(\F,\operatorname{cof}\F,\det\F)$. Polyconvexity implies rank-one convexity and
hence ellipticity of the tangent, and together with suitable coercivity and growth
conditions it guarantees the existence of minimizers~\cite{Ball1977}.
Following~\cite{Klein2022,Linden2023}, the network receives the feature vector
\begin{equation}
  \bm y=\bigl[\,I_1,\ I_2,\ J,\ -\ln J\,\bigr],\qquad
  I_1=\|\F\|^2,\quad I_2=\|\operatorname{cof}\F\|^2 .
\end{equation}
$I_1$ is convex in $\F$, $I_2$ is convex in $\operatorname{cof}\F$, and both $J$
and $-\ln J$ are convex in $J$. An input-convex network~\cite{Amos2017} with
non-negative weights $W^+=\mathrm{softplus}(W)$ and convex, non-decreasing
softplus activations is convex and non-decreasing in $\bm y$, so the composite
energy is polyconvex. The particle volume fraction $\vfrac$ enters only through
an unconstrained sub-network that shifts the hidden-layer biases. This makes
the network \emph{partially} input-convex: convex in the kinematic inputs for
every fixed $\vfrac$. The network uses three hidden layers of width 32 and has
5793 trainable parameters (2561 without the context sub-network, as in Phase~1). Following~\cite{Linden2023}, the raw output is normalized so that the reference
state is energy- and stress-free:
\begin{equation}
  \widehat\Psi(\F,\vfrac)=\Psi_{\mathrm{raw}}(\F,\vfrac)-\Psi_{\mathrm{raw}}(\Iten,\vfrac)
  -p_I(\vfrac)\,(J-1),\qquad
  p_I=\tfrac13\operatorname{tr}\Pstress_{\mathrm{raw}}(\Iten,\vfrac).
  \label{eq:norm}
\end{equation}
The correction depends on $\F$ only through $J$, so it is objective and affine in
a polyconvex argument. The Phase~1 network instead used the correction
$\Pstress_{\mathrm{raw}}(\Iten):(\F-\Iten)$, which is \emph{not} objective,
because $\operatorname{tr}(\Rten-\Iten)\neq0$ for a rotation $\Rten$. The
consequence is quantified in Test~V0d (\cref{sec:v0}).

\section{Periodic RVE with an exactly rigid particle}

The RVE is the unit cube $\Omega=[0,1]^3$ with a centred spherical cavity of
radius $r=(3\vfrac/4\pi)^{1/3}$, $\vfrac\in\{0.02,0.05,0.10\}$
(\cref{fig:rve}).

\begin{figure}[htbp]
\centering
\begin{tikzpicture}[scale=2.1, line join=round]
  % (a) cube with sphere
  \begin{scope}
    \def\d{0.45}
    \fill[tubafblue!8] (0,0)--(1,0)--(1,1)--(0,1)--cycle;
    \fill[tubafblue!14] (1,0)--(1+\d,\d)--(1+\d,1+\d)--(1,1)--cycle;
    \fill[tubafblue!4] (0,1)--(1,1)--(1+\d,1+\d)--(\d,1+\d)--cycle;
    \draw[thick] (0,0)--(1,0)--(1,1)--(0,1)--cycle;
    \draw[thick] (1,0)--(1+\d,\d)--(1+\d,1+\d)--(1,1);
    \draw[thick] (0,1)--(\d,1+\d)--(1+\d,1+\d);
    \draw[dashed,black!50] (0,0)--(\d,\d)--(1+\d,\d); \draw[dashed,black!50] (\d,\d)--(\d,1+\d);
    \shade[ball color=black!35] (0.5+\d/2,0.5+\d/2) circle (0.27);
    \node[font=\small] at (0.72,-0.18) {(a) unit-cell RVE, $\Omega=[0,1]^3$};
    \draw[{Latex}-{Latex},red!70!black,thick] (-0.12,0.5) to[bend left=50] node[above,font=\scriptsize,pos=0.5]{} (1.0+\d+0.12,0.5+\d/2);
    \node[red!70!black,font=\scriptsize,align=center] at (0.72,1.72) {opposite faces identified:\\ $\bm w(\bm X^+)=\bm w(\bm X^-)$};
  \end{scope}
  % (b) cross-section
  \begin{scope}[xshift=2.35cm]
    \fill[tubafblue!8] (0,0) rectangle (1.4,1.4);
    \draw[thick] (0,0) rectangle (1.4,1.4);
    \fill[black!35] (0.7,0.7) circle (0.36);
    \draw[thick] (0.7,0.7) circle (0.36);
    \draw[->,thick] (0.7,0.7) -- ++(25:0.32) node[pos=1.15,font=\scriptsize]{$\Rten$};
    \draw[thick,->] (0.7,0.7) ++(160:0.2) arc (160:250:0.2);
    \node[font=\scriptsize,align=center] at (0.7,1.2) {PDMS: $\mu=1$, $\kappa=10^4$\\ mixed $P_2/P_1$};
    \node[font=\scriptsize,align=center,white] at (0.7,0.58) {rigid};
    \node[font=\scriptsize,align=center] at (0.7,0.17) {$\F=\Fbar+\nabla\bm w$};
    \node[font=\small] at (0.7,-0.18) {(b) cross-section};
  \end{scope}
\end{tikzpicture}
\caption[Microscale problem]{Microscale problem. (a) Periodic unit cell with one centred spherical
particle; opposite faces share their fluctuation degrees of freedom. (b) The
particle surface follows the exact rigid map \cref{eq:rigid}; the free rotation
$\Rten$ is found from moment equilibrium. Volume fractions
$\vfrac=0.02,0.05,0.10$ correspond to radii $0.168$, $0.229$ and $0.288$.
Schematic, not to scale.}
\label{fig:rve}
\end{figure}

The microscopic deformation splits into the prescribed
macroscopic part and a periodic fluctuation,
\begin{equation}
  \bm\varphi(\bm X)=\Fbar\bm X+\bm w(\bm X),\qquad
  \F=\Fbar+\nabla\bm w,\qquad \bm w(\bm X^+)=\bm w(\bm X^-),
  \label{eq:kin}
\end{equation}
so that tractions are anti-periodic in the weak sense. The near-incompressible
PDMS matrix is described by the mixed potential
\begin{equation}
  \Pi(\F,p)=\tfrac{\mu}{2}\bigl(J^{-2/3}I_1-3\bigr)+p\,(J-1)-\frac{p^2}{2\kappa},
  \qquad \mu=1,\ \kappa=10^4 ,
  \label{eq:mixed}
\end{equation}
whose stationarity in $p$ gives $p=\kappa(J-1)$. The particle is not modelled as
a stiff elastic phase. Instead, the cavity surface undergoes the exact finite
rigid motion
\begin{equation}
  \bm u_{\mathrm p}(\bm X)=(\Fbar-\Iten)\bm X_c+(\Rten-\Iten)(\bm X-\bm X_c),\qquad
  \Rten\in SO(3),
  \label{eq:rigid}
\end{equation}
where the centre $\bm X_c$ follows the macroscopic map; a common translation is a
gauge mode of the periodic problem. The rotation is free and is determined by a
vanishing resultant moment on the particle. The Hill--Mandel condition demands
equality of macroscopic and averaged microscopic power~\cite{Hill1963}; at finite
strain the work-conjugate pair is $(\Pbar,\Fbar)$~\cite{Hill1972}:
\begin{equation}
  \Pbar:\dot{\Fbar}=\frac{1}{|\Omega|}\int_{\Omega_{\mathrm m}}\Pstress:\dot{\F}\dd V
  \quad\Longrightarrow\quad
  \Pbar=\frac{1}{|\Omega|}\Bigl(\int_{\Omega_{\mathrm m}}\Pstress\dd V
  +\textstyle\sum_{iJ}(\bm r\cdot\bm b_{iJ})\,\bm e_i\otimes\bm E_J\Bigr),
  \label{eq:hm}
\end{equation}
where $\bm r$ is the assembled nodal residual of the matrix, which contains the
particle reactions, and $\bm b_{iJ}$ is the nodal lift of the particle trace
$-(\bm e_i\otimes\bm E_J)(\bm X-\bm X_c)$ produced by a unit increment of
$\bar F_{iJ}$. The second term is therefore the work of the particle reactions in
the macroscopic modes. Integrating the stress over the
deformable matrix alone would omit this work, and the result would not be
conjugate to $\Fbar$.

\section{Consistent homogenized tangent}

Let $\bm z$ collect the free fluctuation, pressure and rotation unknowns of the
discretized RVE and let $\bm H$ be the Hessian of its discrete energy with
respect to $(\Fbar,\bm z)$. Eliminating the equilibrated internal unknowns gives
the exact derivative of the homogenized stress as a Schur
complement~\cite{Miehe2002},
\begin{equation}
  \Abar=\frac{1}{|\Omega|}\Bigl(\bm H_{\bar F\bar F}-\bm H_{\bar Fz}\,
  \bm H_{zz}^{-1}\,\bm H_{z\bar F}\Bigr).
  \label{eq:schur}
\end{equation}
It is a Hessian of the relaxed energy and must therefore be major-symmetric.
Finite perturbations of $\Fbar$ serve only as verification.

\section{Small-strain reference and scale separation}\label{sec:scales}

At $\Fbar=\Iten$ the homogenized tangent of an isotropic composite reduces to an
effective bulk modulus $K_{\mathrm{eff}}$ and shear modulus $\mu_{\mathrm{eff}}$.
For rigid spheres in a matrix $(\kappa,\mu)$, the Mori--Tanaka
estimate~\cite{MoriTanaka1973}, in the form given by Benveniste~\cite{Benveniste1987},
gives
\begin{equation}
  K_{\mathrm{MT}}=\kappa+\frac{\vfrac\,(\kappa+\tfrac43\mu)}{1-\vfrac},\qquad
  \mu_{\mathrm{MT}}=\mu+\frac{\vfrac\,(\mu+\eta)}{1-\vfrac},\qquad
  \eta=\frac{\mu\,(9\kappa+8\mu)}{6(\kappa+2\mu)} .
  \label{eq:mt}
\end{equation}
For a stiffer inclusion this estimate coincides with the Hashin--Shtrikman lower
bound~\cite{HashinShtrikman1963}, and it serves as an independent check of the RVE.

Of equal importance for the following chapters is the separation of scales
inside the energy. For small strains $\bm\varepsilon$,
\begin{equation}
  \Psibar\approx\tfrac12K_{\mathrm{eff}}(\operatorname{tr}\bm\varepsilon)^2
  +\mu_{\mathrm{eff}}\,\|\operatorname{dev}\bm\varepsilon\|^2,\qquad
  K_{\mathrm{eff}}/\mu_{\mathrm{eff}}\approx(0.87\text{--}0.97)\times10^4 ,
  \label{eq:split}
\end{equation}
where the ratio follows from \cref{eq:mt} for $\vfrac=0.10$ to $0.02$.
Any measure that adds volumetric and deviatoric contributions, such as a
training loss or an error norm normalized by the total energy or stress, is
dominated by the volumetric part unless the volumetric strains are smaller than
the deviatoric ones by a factor of about $\sqrt{\mu/K}\approx10^{-2}$.
\Cref{ch:assessment} shows that this is precisely what limits the present
surrogate.
